% Setup sheet for participants of "From Paper to Proof".
% Build: Overleaf with the pdfLaTeX compiler (or pdflatex)
\documentclass[11pt]{article}
\usepackage[a4paper, margin=1.75cm]{geometry}
\usepackage[T1]{fontenc}
\usepackage[utf8]{inputenc}
\usepackage[sfdefault, lining, semibold, scale=0.92]{FiraSans}
\usepackage[scale=0.92]{FiraMono}
\usepackage{amssymb}
\usepackage{newunicodechar}
\newunicodechar{∀}{\ensuremath{\forall}}
\usepackage{xcolor}
\definecolor{leanblue}{HTML}{2F5DA8}
\definecolor{muted}{HTML}{6B7280}
\usepackage{enumitem}
\setlist{itemsep=2pt, topsep=4pt}
\usepackage[hidelinks]{hyperref}
\usepackage{titlesec}
\titleformat{\section}{\large\bfseries\color{leanblue}}{\thesection}{0.6em}{}
\titlespacing*{\section}{0pt}{1.0em}{0.3em}
\setlength{\parindent}{0pt}
\setlength{\parskip}{0.4em}
\pagestyle{empty}
\sloppy

\begin{document}

{\LARGE\bfseries From Paper to Proof\par}
\vspace{0.2em}
{\large Getting Lean running \textperiodcentered{} BGSMath, UPC, November 2026\par}
\vspace{0.6em}

You don't need a laptop in class: the sessions are demonstrations, and you redo them at home.
For the first two sessions a web browser is all you need. For sessions 3 and 4 it's worth having Lean on your own laptop, and this sheet takes you through it. If anything here doesn't work, don't fight it: use the browser, and ask me at the next session.

\section{In the browser (sessions 1 and 2)}

\begin{enumerate}
\item Go to \href{https://live.lean-lang.org}{\texttt{live.lean-lang.org}}. There's nothing to install.
\item Open a course file, say \texttt{Exercises1.lean}, in the \texttt{BGSMath} folder of \href{https://github.com/mkaratarakis/from-paper-to-proof}{\texttt{github.com/mkaratarakis/from-paper-to-proof}}. Click \emph{Raw}, then copy all of it.
\item Back in the browser, delete whatever is in the editor and paste the file in.
\item The first line, \texttt{import Mathlib}, takes a minute or two to load the first time. That's normal.
\item The panel on the right is the \emph{infoview}. Click on any line and it shows you what Lean knows there.
\end{enumerate}
The page doesn't save your work, so copy it back into a file before you close it.

\section{On your laptop (sessions 3 and 4)}

Set aside about half an hour, a good internet connection, and a few gigabytes of disk space.
These are the steps from \href{https://lean-lang.org/install/}{\texttt{lean-lang.org/install}}, which also has screenshots.

\begin{enumerate}
\item Install \textbf{VS Code} from \href{https://code.visualstudio.com}{\texttt{code.visualstudio.com}}. It's a free text editor; Lean lives inside it.
\item In VS Code, open the Extensions panel (Ctrl+Shift+X, or Cmd+Shift+X on a Mac), search for \texttt{lean4} and install \textbf{Lean 4}. A setup guide opens; follow it, and it installs Lean for you.
\item Create a project with Mathlib: click the \texttt{∀} symbol at the top right of VS Code, then \emph{New Project…} and \emph{Create Project Using Mathlib…}. Pick a folder and a name, say \texttt{bgsmath-lean}. The first time, Mathlib is downloaded, which takes a while.
\item Get the course files from the same address as above: \emph{Code} → \emph{Download ZIP}. Copy the \texttt{.lean} files from its \texttt{BGSMath} folder into your project folder.
\item \textbf{Try it.} Open \texttt{Test.lean} from your project folder and click on the line \texttt{\#eval 2 + 2}. If the infoview shows \texttt{4} and nothing is underlined in red, you're done.
\end{enumerate}

\textbf{One thing that can go wrong:} Lean starts ``building'' thousands of files. Stop it: close VS Code, open a terminal in the project folder, run \texttt{lake exe cache get}, and open VS Code again.

\section{An AI agent (session 3, optional)}

\begin{itemize}
\item In session 3 I'll use a coding agent, Claude Code, to write Lean proofs. Other agents work the same way.
\item Install instructions: \href{https://code.claude.com/docs}{\texttt{code.claude.com/docs}}. You need an account, and most plans are paid, so this is optional.
\item To let the agent see what you see in the infoview, add \textbf{lean-lsp-mcp}, following \href{https://github.com/oOo0oOo/lean-lsp-mcp}{\texttt{github.com/oOo0oOo/lean-lsp-mcp}}.
\item No agent? Every exercise can be done by hand.
\end{itemize}

\section{If you're stuck}

\begin{itemize}
\item Use the browser for now; it's enough for sessions 1 and 2.
\item Ask at the start of the next session, or on the Lean Zulip chat, in the stream \emph{\#new members}: \href{https://leanprover.zulipchat.com}{\texttt{leanprover.zulipchat.com}}. Questions of every level are welcome there, really.
\end{itemize}

\end{document}
